% chapters/02_sections/dram_and_hbm.tex
\subsection{DRAM Fundamentals and High Bandwidth Memory}\label{sec:dram-hbm}

Understanding the physical memory technology that backs GPU global memory is
important for reasoning about bandwidth limits, access latency, and power
behavior. This subsection summarizes the relevant DRAM fundamentals and then
describes the High Bandwidth Memory (HBM) family used in modern GPU
accelerators, with emphasis on the HBM2e variant deployed in the
A100-SXM4-80GB.

\subsubsection{The DRAM Cell}\label{sec:dram-cell}

Dynamic Random Access Memory stores each bit as charge on a single capacitor,
gated by a single access transistor. This \emph{1T1C} (one-transistor,
one-capacitor) design is the dominant commodity DRAM cell and has been since
the early 1970s~\cite{li2018highspeed-dram}. In steady state, a charged
capacitor represents a logical~1 and a discharged capacitor a logical~0 (or
vice versa, depending on convention). The access transistor connects the
capacitor to a shared \emph{bitline} (column wire) when the transistor's gate
is driven high by the \emph{wordline} (row wire).

Because the capacitor is small (on the order of femtofarads), two consequences
follow. First, the stored charge leaks over time, so every cell must be
periodically \emph{refreshed}---read and rewritten---to preserve its content;
this is the ``dynamic'' in DRAM. Typical refresh intervals are
\SI{32}{\milli\second} to \SI{64}{\milli\second} for the entire
array~\cite{li2018highspeed-dram}. Second, reading a cell is
\emph{destructive}: the act of sharing charge between the cell capacitor and
the bitline disturbs the stored voltage, so the sense amplifier must both
detect the charge and immediately restore it.

\paragraph{Connection to energy and data-dependent power.}
Switching a bit between 0 and 1 requires charging or discharging the cell
capacitor and toggling the I/O drivers on the data bus. The classical
switching-energy model gives a per-transition energy proportional to
$\tfrac{1}{2}CV^2$ (charging/discharging a capacitance $C$ through a voltage
swing $V$). When many bits toggle simultaneously---as happens when writing
data with high Hamming distance from the previous content---aggregate
switching energy increases. Conversely, data patterns with low bit-toggle
rates (e.g.\ writing a matrix of zeros, or repeated values) cause fewer
transitions and consume less I/O power.

This data-dependent power behavior is not merely theoretical. HBM2e devices
implement \emph{Data Bus Inversion} (DBI), a per-byte encoding that inverts
the data on the bus whenever more than half the bits would toggle, thereby
capping the worst-case toggle rate at 50\%~\cite{micron2021hbm2e}. The
mechanism is a direct hardware response to switching-energy concerns.
On the GPU side, the same physics underlies the data-pattern sensitivity
explored experimentally in \cref{sec:energy-gemm-switching-todo}: if the
memory subsystem throttles power under sustained high-toggle workloads, the
effect can propagate into observable performance variation.

\subsubsection{Row, Column, and Bank Organization}\label{sec:dram-organization}

Internally, DRAM cells are organized in a two-dimensional array of rows and
columns, as illustrated schematically in Figure~2
of Li et al.~\cite{li2018highspeed-dram}. A read proceeds in three phases:

\begin{enumerate}
  \item \textbf{Row activation} (\emph{activate}, timing parameter
    $t_\text{RCD}$): the wordline for the target row is raised, connecting
    all capacitors in that row to their respective bitlines. The sense
    amplifiers detect and amplify the charge, effectively copying the entire
    row into the \emph{row buffer} (sense-amplifier latch array). This is the
    most time-consuming step and also restores the cells that were
    destructively read.
  \item \textbf{Column access} (timing parameter $t_\text{CL}$): a column
    address selects a subset of the row buffer and drives it onto the I/O
    pins. Because the full row is already latched, subsequent reads to
    different columns in the same row (``row-buffer hits'') require only this
    shorter column-access latency.
  \item \textbf{Precharge} (timing parameter $t_\text{RP}$): before a
    different row in the same bank can be activated, the bitlines must be
    precharged to a reference voltage. This closes the current row.
\end{enumerate}

The distinction between a \emph{row-buffer hit} (column access only,
cost $\approx t_\text{CL}$) and a \emph{row-buffer miss} requiring a full
row switch (cost $\approx t_\text{RP} + t_\text{RCD} + t_\text{CL}$) is
central to DRAM performance; the existing discussion in
\cref{sec:hbm-row-column-path} develops the quantitative implications for
kernel design.

To increase parallelism, the array is partitioned into \emph{banks} that can
be activated independently. Multiple outstanding requests to different banks
can overlap, which is how high-bandwidth memory technologies sustain
throughput despite individual access latencies in the tens of
nanoseconds~\cite{li2018highspeed-dram}.

\paragraph{Evolution of DRAM bandwidth.}
Historically, increasing DRAM bandwidth has been far easier than decreasing
latency. The core timing parameters ($t_\text{RCD}$, $t_\text{RP}$,
$t_\text{CL}$) are constrained by physical properties of the storage
capacitors and bitlines; they have decreased only modestly over two decades.
Bandwidth improvements have instead come from widening the data path (the
\emph{prefetch buffer} between core and I/O), increasing pin data rates, and
adding more channels and banks. This observation---articulated clearly
by Li et al.~\cite{li2018highspeed-dram}---explains why the HBM approach (massively
wider buses rather than faster pins) is so effective.

\paragraph{Where bandwidth is lost in practice.}
Even with wide buses and many banks, achieved DRAM bandwidth is always
below the theoretical peak. Eyerman et al.~\cite{eyerman2022dram-stacks}
introduce \emph{bandwidth stacks} and \emph{latency stacks} as a
decomposition methodology: for each memory-controller cycle that does
\emph{not} transfer useful data, the lost bandwidth is attributed to a
specific cause---row-buffer misses (precharge + activate overhead), DRAM
refresh, bus turnaround between reads and writes, or queuing delays from
timing constraints. Their analysis shows that row conflicts and refresh
overhead are often the dominant bandwidth losses, while write-to-read bus
turnaround becomes significant in mixed read/write traffic.

This decomposition connects directly to GPU kernel design. Coalesced
warp accesses maximise row-buffer hit rates, reducing the
precharge/activate fraction. Tiled algorithms that operate on contiguous
blocks of memory reduce the frequency of row switches within a bank.
And because HBM refresh is non-negotiable (the 1T1C cell requires it),
the refresh overhead sets a hard floor on bandwidth loss that no access
pattern can avoid. The practical implication is that an HBM channel
delivering 90--95\% of its peak bandwidth under a perfectly sequential
stream is already close to optimal; the kernel optimiser's job is to
keep real workloads as close to that ceiling as
possible~\cite{eyerman2022dram-stacks,li2018highspeed-dram}.

\subsubsection{High Bandwidth Memory}\label{sec:hbm-architecture}

JEDEC's High Bandwidth Memory standard addresses the bandwidth gap by using
3D integration to package a stack of DRAM dies together with a base logic
die~\cite{jedec2013hbm,jedec2016hbm2,jedec2020hbm2e}. The stack communicates
internally via \emph{through-silicon vias} (TSVs)---vertical conductors
etched through the silicon substrate---which enable a very wide parallel bus
within the stack (1024 data bits across 8 independent channels of 128 bits
each). Jun et al.~\cite{jun2017hbm} describe the core engineering
challenges of this stacking approach---TSV reliability, thermal dissipation
through the stack, and high-throughput testing of individual dies before
assembly---and provide a generational comparison of the HBM family from the
manufacturer's perspective.

The entire HBM stack is mounted on a \emph{silicon interposer} alongside the
GPU die, forming a \emph{system-in-package} (SiP). This 2.5D packaging
approach avoids the signal-routing constraints of a conventional PCB, where
trace count is limited to roughly 384 data lines per
SoC~\cite{micron2021hbm2e}. Because the interposer traces are much shorter
and denser than PCB traces, HBM achieves its bandwidth at comparatively low
per-pin data rates (around \SI{1}{\giga bps} for HBM1, \SI{2}{\giga bps}
for HBM2, and \SI{3.2}{\giga bps} for HBM2e), which in turn reduces I/O
power substantially compared with high-speed serial interfaces like GDDR6
that must drive pins at \SI{16}{\giga bps} or
more~\cite{li2018highspeed-dram,micron2021hbm2e}.

\paragraph{Memory dies and channels.}
The DRAM layers in an HBM stack are called \emph{memory dies} (or
\emph{core dies}). Each memory die is a complete DRAM chip in its own right,
containing billions of 1T1C storage cells organized into banks, rows, and
columns. The dies are manufactured and tested individually before being
stacked; Micron ships them as \emph{Known Good Stacked Die}
(KGSD)~\cite{micron2021hbm2e}.

A \emph{channel} is a fully independent memory interface with its own clock,
command/address bus, and 128-bit data path (204 signals in
total)~\cite{micron2021hbm2e}. Because each channel has its own controller
state, eight outstanding memory transactions---one per channel---can proceed
in parallel. Each of the 8 channels is further subdivided into two
\emph{pseudo channels} (PCs) that share clock and command inputs but have
independent 64-bit data paths, effectively providing 16 semi-independent
access points per device~\cite{micron2021hbm2e}. This high degree of
internal parallelism is what allows HBM to sustain throughput even when
individual DRAM timing parameters ($t_\text{RCD}$, $t_\text{RP}$,
$t_\text{CL}$) are comparable to conventional GDDR.

\paragraph{Banks.}
Within each pseudo channel the DRAM array is partitioned into 16
\emph{banks} (B0--B15) that can be activated independently.
A bank is the smallest unit that can hold an open row: each bank has its
own row buffer (sense-amplifier latch), so accesses to different banks
can overlap without incurring row-switch penalties. In the 4-Hi configuration each channel has
16 banks; in 8-Hi, two dies contribute to the same channel, yielding
$2 \times 16 = 32$ banks per channel (distinguished by a \emph{stack ID}
bit that selects the upper or lower memory die)~\cite{micron2021hbm2e}.

\paragraph{Channel-to-die mapping.}
Channels are distributed across the memory dies in pairs.
In a 4-Hi stack, each of the four core dies carries two channels (e.g.\ the
bottom die carries Channels~A and~C, the next die Channels~B and~D, and so
on). In an 8-Hi stack the same channel pair is replicated across two dies,
doubling the capacity per channel while preserving the same 8-channel
topology. \Cref{fig:hbm2e-stack-side} illustrates this mapping for both
configurations: blue layers are memory core dies, the yellow layer is the
interface (base) die, and the red layer in the 8-Hi diagram marks a stack
position that may be disabled for yield.
\Cref{fig:hbm2e-die-topdown} shows the internal layout of a single memory
die, where each pseudo channel's 16 banks are visible as individually
addressable units.

\begin{figure}[H]
  \centering
  \includegraphics[width=0.92\textwidth]{HBM2EStackSideView}
  \caption{HBM2e stack side view: 4-Hi (left) vs.\ 8-Hi (right).
  Channel pairs are labelled per die.
  After~\cite{micron2021hbm2e}, Fig.~4.}
  \label{fig:hbm2e-stack-side}
\end{figure}

\begin{figure}[H]
  \centering
  \includegraphics[width=0.65\textwidth]{HBM2EDieTopDown}
  \caption{Top-down view of a single HBM2e memory die with two pseudo
  channels (A, C) and their 16~banks each (B0--B15).
  After~\cite{micron2021hbm2e}, Fig.~4.}
  \label{fig:hbm2e-die-topdown}
\end{figure}

\paragraph{HBM generation roadmap.}
\Cref{tab:hbm-generations} summarises the key parameters across HBM
generations relevant to datacenter GPUs. Through HBM2e the 1024-bit bus
and 8-channel topology are preserved; bandwidth scaling comes almost
entirely from increasing the per-pin data rate, while capacity scaling
comes from stacking more dies (4-Hi $\to$ 8-Hi $\to$ 12-Hi) and from
denser per-die DRAM arrays. HBM3/3e doubles the channel count to~16
(each now 64-bit) while retaining the 1024-bit aggregate
width~\cite{jun2017hbm,jedec2020hbm2e}.

\begin{table}[H]
\centering
\caption{HBM generation comparison for datacenter-relevant configurations.
  Data compiled from JEDEC standards~\cite{jedec2013hbm,jedec2016hbm2,jedec2020hbm2e}
  and Jun et al.~\cite{jun2017hbm}.}
\label{tab:hbm-generations}
\begin{tabular}{l c c c c c}
\toprule
& HBM & HBM2 & HBM2e & HBM3 & HBM3e \\
\midrule
JEDEC adoption         & 2013    & 2016    & 2018     & 2022    & 2023 \\
Pin data rate          & 1 Gbps  & 2 Gbps  & 3.2 Gbps & 6.4 Gbps & 9.6 Gbps \\
Bus width (bits)       & 1024    & 1024    & 1024     & 1024    & 1024 \\
Channels per stack     & 8       & 8       & 8        & 16      & 16 \\
Max stack height       & 4-Hi    & 8-Hi    & 8-Hi     & 8/12-Hi & 12-Hi \\
Max capacity / stack   & 4 GB    & 8 GB    & 16 GB    & 24 GB   & 36 GB \\
BW per stack           & 128 GB/s & 256 GB/s & 410 GB/s & 819 GB/s & 1229 GB/s \\
\bottomrule
\end{tabular}
\end{table}

\subsubsection{HBM2e vs.\ GDDR6(X) and the A100 Memory Subsystem}
\label{sec:hbm-generations}

HBM2e is often positioned as an alternative to the well-established
GDDR6 and GDDR6X SGRAM technologies, which also provide high memory
bandwidth but through a fundamentally different
trade-off~\cite{micron2021hbm2e}. Where HBM2e achieves bandwidth through a
$32\times$ wider I/O bus at modest pin speeds (\SI{3.2}{Gbps}), GDDR6(X)
compensates for its narrow 32-bit bus by driving pins at
\SI{16}--\SI{24}{Gbps}, which demands power-hungry techniques such as
on-die signal termination, delay-locked loops, and clock-data
recovery~\cite{micron2021hbm2e}. \Cref{tab:hbm-vs-gddr} contrasts the
key parameters.

The V100 (Volta, 2017) was the first NVIDIA GPU accelerator to deploy HBM2
at scale, establishing the pattern of in-package high-bandwidth memory that
the A100 continues.

\begin{table}[H]
\centering
\caption{HBM2e DRAM vs.\ GDDR6(X) SGRAM.
  Data from~\cite{micron2021hbm2e}.}
\label{tab:hbm-vs-gddr}
\begin{tabularx}{\textwidth}{l r r}
\toprule
Parameter & HBM2e & GDDR6 / GDDR6X \\
\midrule
Capacity per device    & 8--\SI{16}{GB}       & 1--\SI{2}{GB} \\
Channels per device    & 8                     & 2 \\
Devices per system     & 4--6                  & up to 24 \\
Max pin data rate      & \SI{3.2}{Gbps}        & 16 / \SI{24}{Gbps} \\
I/O width (bits)       & 1024                  & 32 \\
BW per device          & \SI{410}{GB/s}        & \SI{64}{GB/s} \\
Total system BW        & up to \SI{2450}{GB/s} & up to 768 / \SI{1152}{GB/s} \\
Package                & KGSD (stacked die)    & BGA-180 \\
\bottomrule
\end{tabularx}
\end{table}

Several observations from this comparison connect directly to the performance
modeling in later chapters:

\begin{enumerate}
  \item \textbf{Latency is comparable.} The core DRAM timing parameters
    ($t_\text{RCD}$, $t_\text{RP}$, $t_\text{CL}$) are in the same
    $\sim$\SI{14}{\nano\second} range for both HBM2e and GDDR6. High
    bandwidth does not imply low latency; latency hiding requires sufficient
    parallelism (banks, channels, and outstanding
    requests)~\cite{li2018highspeed-dram}.
  \item \textbf{HBM2 to HBM2e is a pin-speed upgrade.} The physical
    interface (interposer, TSV pitch, channel count) is unchanged between
    HBM2 and HBM2e. The bandwidth increase from \SI{256}{GB/s} to
    \SI{410}{GB/s} per stack comes from raising the per-pin data rate from
    \SI{2}{Gbps} to \SI{3.2}{Gbps}~\cite{micron2021hbm2e}. HBM2e also
    doubles the maximum stack capacity from \SI{8}{GB} to \SI{16}{GB}
    (eight-high configuration with denser dies), which is why the
    A100-SXM4-80GB can offer \SI{80}{GB} from five stacks.
  \item \textbf{GPU deployment history.} The V100 (Volta, 2017) used four
    stacks of HBM2 providing \SI{900}{GB/s} aggregate bandwidth and
    \SI{16}{GB} or \SI{32}{GB} capacity. Both A100 variants use five active
    HBM stacks out of six physical stacks on the package (one disabled for
    yield), giving $5 \times 8 = 40$ independent channels in both cases.
    The A100-SXM4-40GB uses HBM2 in a \emph{4-Hi} (four-high) stack
    configuration: \SI{1}{GB} per die $\times$ 4~dies $=$ \SI{8}{GB} per
    stack, for \SI{40}{GB} total at \SI{1555}{GB/s}. The A100-SXM4-80GB
    upgrades to HBM2e in an \emph{8-Hi} (eight-high) configuration:
    \SI{2}{GB} per die $\times$ 8~dies $=$ \SI{16}{GB} per stack, for
    \SI{80}{GB} total at \SI{2039}{GB/s}~\cite{nvidia2020a100,micron2021hbm2e}.
    The stacks are physically twice as tall but share the same channel
    topology; the bandwidth increase comes from the higher per-pin data rate
    of HBM2e (\SI{3.2}{Gbps} vs.\ $\sim$\SI{2.4}{Gbps} for the HBM2
    configuration used in A100-40GB). The quantitative impact of this
    bandwidth difference on roofline modeling is developed in
    \cref{sec:a100-40-80}.
  \item \textbf{Hopper generation.} The H100-SXM5-80GB (2022) moved to
    HBM3 with six stacks of \SI{13.3}{GB} each, providing \SI{80}{GB} at
    \SI{3350}{GB/s}~\cite{nvidia2022h100}. The H200 (2024) shares the
    same GH100 die but upgrades to HBM3e: six stacks totalling
    \SI{141}{GB} at \SI{4800}{GB/s}. This pattern---same compute die,
    upgraded memory---mirrors the A100-40GB$\to$80GB transition and
    illustrates that memory technology often advances faster than a
    GPU generation cycle.
\end{enumerate}
